\section{From local divisibility to a global arithmetic bound}
\label{growth:section}

The preceding section controls each prime separately. We now add those
contributions and show that making the coefficients integral costs less
than the real decay can afford. This is Fauzan's normalization strategy,
with the revised local estimates and cutoff used in Firsching's
formalization. In particular, the outer range uses the full rank loss;
we do not discard the correction discussed in Lai's write-up.
\cite[Section 5]{fauzan}
\cite[\texttt{Norm.lean}, \texttt{Growth.lean}]{firsching}

Throughout this section,
\[
 K=40n,\qquad N=3n,\qquad h=37n,\qquad
 \alpha=\frac3{40},\qquad \lambda=\frac{37}{40},\qquad M=400.
\]
The cutoff \(M\) is fixed while \(n\) tends to infinity. This matters:
an error bounded in terms of \(M\) is independent of \(n\).

The profiles below summarize the local valuation bounds. There are two
passages from sums to integrals: averaging residue classes \(z=c/p\),
then summing the resulting costs over primes \(p\). The variable
\(x=K/p\) keeps track of the number of poles in each class.

\subsection{The multiplier and the primes that occur}

Recall that \(F_K=S_K\Delta_K\), where
\[
 S_K=\frac{(K!)^{2h}4^{h-1}}
 {(N!)^{12h}\prod_{i=1}^{h-1}((2i)!)^2}.
\]
Write \(v_p^G(F_K)\) for the minimum valuation of its nonzero
coefficients. For all sufficiently large \(n\), define the following
integer lower bound for \(v_p^G(F_K)\):
\begin{equation}\label{growth:Lp}
 L_p=
 \begin{cases}
 -6h\lfloor\log_p(5K)\rfloor-hv_p(24),&p\le K/400,\\
 v_p(S_K)+\lfloor 2W_p^{\rm in}\rfloor,&K/400<p\le K/3,\\
 v_p(S_K)+\lfloor W_p^{\rm out}\rfloor,&K/3<p\le 2h.
 \end{cases}
\end{equation}
Here \(W_p^{\rm in}\) is the sum of the capped inner weights constructed
in the preceding section, and \(W_p^{\rm out}\) is twice the sum of the
outer weights minus the full rank loss. These are precisely the
quantities denoted \texttt{WI} and \texttt{WO} in
\cite[\texttt{Norm.lean}]{firsching}. The local results give
\[
 v_p^G(F_K)\ge L_p.
\]
The hypotheses of those results hold uniformly in the last two ranges
once \(n\) is sufficiently large. Indeed \(p>K/400=n/10\), so
\(p^2>n^2/100>104n+6\) for \(n\ge12001\). This implies the pole-separation
and degree conditions; it also implies \(p\ge7\). The bounded reserve of
rows in the zero class fits into \(h=37n\).

No prime larger than \(2h\) is needed. At such a prime all factorial
arguments in \(S_K\) are smaller than \(p\), so \(v_p(S_K)=0\).
Every pole index \(j\le K\) is also smaller than \(p\). In the outer
local bound there are therefore no large-pole rows, every outer weight
is zero, and the rank loss is zero. Consequently
\begin{equation}\label{growth:big}
 v_p^G(F_K)\ge0\qquad(p>2h).
\end{equation}
This is the argument recorded as \texttt{GV\_F\_big} in
\cite[\texttt{Norm.lean}]{firsching}.

Set
\begin{equation}\label{growth:normalizer}
 m_K=\prod_{\substack{p\le2h\\p\ {\rm prime}}}p^{-L_p}.
\end{equation}
It is a positive rational number. For each coefficient \(a\) of \(F_K\),
\[
 v_p(m_Ka)\ge0
\]
at every prime, by \eqref{growth:Lp} and \eqref{growth:big}. A reduced
rational number with this property has denominator one: any prime
dividing a denominator would give a negative valuation. Thus
\begin{equation}\label{growth:integrality}
 m_KF_K\in\mathbb Z[X],\qquad
 \log m_K=\sum_{p\le2h}(-L_p)\log p.
\end{equation}
Positive \(L_p\)'s contribute negative terms to this logarithm. They
record common factors of the coefficients that can be divided out.
The calculation controls both clearing denominators and removing such
common factors.

\subsection{Why sums over primes become integrals}

We use the prime number theorem in the standard form
\[
 \vartheta(T):=\sum_{p\le T}\log p=T+o(T).
\]
The following consequence explains the factor \(x^{-3}\) in all the
integrals below. It is the prime-summation step of
\cite[Section 5]{fauzan}; the formalization supplies it in
\cite[\texttt{PrimeSum.lean}, \texttt{PrimeSumIntegral.lean}]{firsching}.

\begin{lemma}\label{growth:prime-sum}
If \(0<a<b\) are fixed and \(E\) is bounded and piecewise continuous on
\([a,b]\), with only finitely many discontinuities, then
\[
 \frac1{K^2}\sum_{K/b<p\le K/a}p\log p\,E(K/p)
 \longrightarrow \int_a^b\frac{E(x)}{x^3}\,dx .
\]
\end{lemma}
\begin{proof}
The prime number theorem implies, for every fixed \(s<t\),
\[
 \frac1K\sum_{sK<p\le tK}\log p\longrightarrow t-s.
\]
Hence the measures
\(K^{-1}\sum_p(\log p)\delta_{p/K}\), restricted to the compact interval
\([1/b,1/a]\), converge against step functions to ordinary length
measure. Approximating a continuous function uniformly by step functions
gives the same assertion for continuous functions. A bounded function
with finitely many discontinuities is treated by isolating those points
in intervals of arbitrarily small total length and approximating on
the complement. The mass of these exceptional intervals has the
corresponding small limiting upper bound, again by the displayed
prime-number-theorem limit.

Apply this to \(t\mapsto tE(1/t)\). The resulting integral is
\(\int_{1/b}^{1/a}tE(1/t)\,dt\). The substitution \(x=1/t\) gives the
asserted integral.
\end{proof}

In particular, a per-prime estimate
\[
 -L_p\le pE(K/p)+O_M(1)
\]
can be summed with no quadratic cost from its error:
\[
 \sum_{p\le2h}O_M(1)\log p=O_M(\vartheta(2h))=O_M(K)=o(K^2).
\]
Values at the finitely many breakpoints of a profile do not affect
the limit. For a fixed breakpoint \(x_0\), the equation \(K/p=x_0\)
holds for at most one prime, whose contribution is \(O(K\log K)\).

\subsection{Small primes}

For \(p\le K/400\), the first line of \eqref{growth:Lp} gives
\[
 \sum_{p\le K/400}(-L_p)\log p
 =
 6h\sum_{p\le K/400}\lfloor\log_p(5K)\rfloor\log p
 +h\sum_{p\le K/400}v_p(24)\log p.
\]
The second sum is at most \(\log24\). In the first sum, primes larger
than \(\sqrt{5K}\) have \(\lfloor\log_p(5K)\rfloor=1\). For the
remaining primes the summand is at most \(\log(5K)\), and there are
at most \(\sqrt{5K}\) such primes. Therefore
\[
 \sum_{p\le K/400}\lfloor\log_p(5K)\rfloor\log p
 \le\vartheta(K/400)+O(\sqrt K\log K).
\]
Since \(h=\lambda K\), the prime number theorem yields
\begin{equation}\label{growth:small}
 \limsup_{K\to\infty}\frac1{K^2}
 \sum_{p\le K/400}(-L_p)\log p
 \le\frac{6\lambda}{400}=\frac{111}{8000}.
\end{equation}

\subsection{The continuous inner profile}

We next turn the residue-class calculation into a fixed function of
\(x=K/p\). This step follows
\cite[\texttt{Growth/ClassSum.lean}, \texttt{ClassData.lean},
\texttt{InnerPrime.lean}, \texttt{TablePrime.lean}]{firsching}.
Put
\[
 x=\frac Kp,\quad u=\lambda x,\quad
 q=\lfloor x\rfloor,\quad q'=\lfloor\alpha x\rfloor,\quad
 f=\{x\},\quad g=\{\alpha x\}.
\]
For \(0<z<1/2\), define
\begin{equation}\label{growth:af}
 a_f(z)=\mathbf1_{z\le f}+\mathbf1_{z\ge1-f}.
\end{equation}
If \(1\le c\le(p-1)/2\), the number of integers \(1\le j\le K\)
with \(j\equiv c\) or \(-c\pmod p\) is exactly
\[
 2q+a_f(c/p).
\]
The analogous count for \(N\) is \(2q'+a_g(c/p)\). Thus the nonzero
class offsets in the inner bound are
\[
 \beta_c=B+6a_g(c/p)-a_f(c/p),\qquad B=12q'-2q-4.
\]

For integers \(k,\beta\), write
\begin{equation}\label{growth:psi}
 \psi(k,\beta)=
 \sum_{i\ge0}\frac{(k-\beta-2i)_+}{2}
 =\frac{\ell(k-\beta-\ell+1)}2,\qquad
 \ell=\max\!\left(0,\left\lfloor\frac{k-\beta+1}{2}\right\rfloor\right).
\end{equation}
Here \(t_+=\max(0,t)\), and the displayed sum has finitely many
nonzero terms. This function measures the total deficit below the
level \(k/2\) in a class with raw weights \(i+\beta/2\).
The water-filling estimate \eqref{arith:inner-comparison} says that,
for every integer
\[
 -2q-6\le k\le8q+20,
\]
\begin{equation}\label{growth:WI}
 2W_p^{\rm in}\ge
 hk-2\sum_{c=1}^{(p-1)/2}\psi(k,\beta_c)-O_M(1).
\end{equation}
The error is bounded by the product of the level range and the reserve,
plus twice the deficit of the initial zero-class rows below the selected
level \(k_*/2\).
Because \(q<400\), both quantities are bounded independently of \(K\).
The total number of zero-class rows need not be bounded; all subsequent
zero-class rows contribute at least \(k_*/2\) each and cause no deficit.
This is the bounded-error form of \texttt{inner\_Lp\_le} in the cited
source.

The factorial normalization contributes another elementary function.
Define
\begin{equation}\label{growth:J}
 J(u)=\sum_{j\ge1}(u-j/2)_+
 =\mu u-\frac{\mu(\mu+1)}4,\qquad \mu=\lfloor2u\rfloor.
\end{equation}
Legendre's formula, and \(p^2>2h\), give
\begin{align*}
 -v_p(S_K)
 &=-2hq+12hq'+2\sum_{i=1}^{h-1}\left\lfloor\frac{2i}{p}\right\rfloor\\
 &\le -2hq+12hq'+2pJ(h/p).
\end{align*}
For the inequality, count the layers
\(\lfloor2i/p\rfloor=\sum_{j\ge1}\mathbf1_{jp\le2i}\).
For fixed \(j\), fewer than or equal to \((h-jp/2)_+\) indices
\(1\le i<h\) occur. Summing over \(j\) gives \(pJ(h/p)\).

Combining this inequality with \eqref{growth:WI}, and allowing one unit
for the floor in \eqref{growth:Lp}, gives
\begin{equation}\label{growth:inner-discrete}
 -L_p\le -2hq+12hq'+2pJ(u)-hk
       +2\sum_{c=1}^{(p-1)/2}\psi(k,\beta_c)+O_M(1).
\end{equation}
The function of \(z\) inside the class sum is constant between the at
most four breakpoints \(f,1-f,g,1-g\). Counting mesh points \(c/p\)
on each of these intervals differs from its length times \(p\)
by a bounded number. Its heights are bounded in terms of \(M\), since
\(q,q'\), \(k\), and \(\beta_c\) are. Hence
\[
 \sum_{c=1}^{(p-1)/2}\psi(k,\beta_c)
 =
 p\int_0^{1/2}\psi(k,B+6a_g(z)-a_f(z))\,dz+O_M(1).
\]
Dividing the principal terms of \eqref{growth:inner-discrete} by \(p\)
therefore produces the explicit profile
\begin{align}
 E_{\rm in}(x,k)
 ={}&-2uq+12uq'+2J(u)-uk \notag\\
 &+2\int_0^{1/2}\psi(k,B+6a_g(z)-a_f(z))\,dz,
 \label{growth:Ein}\\
 -L_p\le{}&pE_{\rm in}(K/p,k)+O_M(1).
 \label{growth:inner-profile-bound}
\end{align}
The integral in \eqref{growth:Ein} is a finite sum of interval lengths,
not an integral requiring numerical quadrature.
The first three terms record the factorial normalization; \(-uk\)
subtracts the common row level, and the integral restores the deficits.
Legendre's floors and the real logarithmic estimate thus examine the
same finite products by counting prime powers and summing logarithms,
respectively.

\subsection{The finite inner range \texorpdfstring{\(3\le x\le20\)}{3 <= x <= 20}}

Choose, on each interval, the admissible integer \(k\) that minimizes
\eqref{growth:Ein}. Denote the resulting piecewise affine function by
\(R_{\rm in}(x)\). The exact calculation gives
\begin{equation}\label{growth:inner-integral}
 \int_3^{20}\frac{R_{\rm in}(x)}{x^3}\,dx
 =\frac{322437603634266857629}{7535670527041937280000}
 =0.042788\ldots.
\end{equation}
This is the rational value in Fauzan's equation (5.18). The revised
formalization derives it from its own class bounds, in
\cite[\texttt{Growth/InnerTable.lean}, \texttt{InnerSum.lean}]{firsching}.

Here are complete instructions for the finite calculation, including
why it is exact. Partition \([3,20]\) at all points in the five families
\begin{equation}\label{growth:grid}
 \frac{\mathbb Z}{2},\qquad
 \frac{20\mathbb Z}{3},\qquad
 \frac{20\mathbb Z}{37},\qquad
 \frac{40\mathbb Z}{37},\qquad
 \frac{40\mathbb Z}{43}.
\end{equation}
The first two families fix the integer parts of \(x,\alpha x\)
and the positions of their fractional parts relative to \(1/2\).
The third fixes \(\mu=\lfloor2\lambda x\rfloor\). The last two
include every crossing between \(f,1-f,g,1-g\). On each open interval
of the partition, these four functions and the endpoints \(0,1/2\)
have a fixed order. Their successive differences are affine functions
of \(x\); the values of \(a_f,a_g\) in each resulting strip are fixed
integers. Formula \eqref{growth:Ein} consequently gives an affine
function \(A_kx+B_k\) for each of the finitely many admissible \(k\)'s.

Take the lower envelope of these lines by inserting their relevant
pairwise intersections. A proposed line is the minimum throughout
an interval if it is no larger than every competitor at both
endpoints: their differences are affine. Finally integrate each
piece using
\begin{equation}\label{growth:affine-integral}
 \int_a^b\frac{Ax+B}{x^3}\,dx
 =A\left(\frac1a-\frac1b\right)
  +\frac B2\left(\frac1{a^2}-\frac1{b^2}\right).
\end{equation}
The resulting table has 125 pieces. Its full rational data are
printed in Appendix~\ref{growth:table-appendix} and supplied in
\texttt{checks/growth\_inner\_table.csv};
the accompanying \texttt{growth\_exact.py} implements the branch
construction and checks all endpoint inequalities with rational
arithmetic. Thus \eqref{growth:inner-integral} is reproducible by
finite exact arithmetic. No grid sampling or floating-point
optimization is used to certify it.

\subsection{The inner tail \texorpdfstring{\(20\le x\le400\)}{20 <= x <= 400}}

For this longer interval, a quadratic upper bound is more convenient
than a large table. We give the calculation, which is the continuous
form of \cite[\texttt{Growth/TailPrime.lean}]{firsching}.
Completing the square in \eqref{growth:psi} shows
\[
 \psi(k,\beta)\le\frac{(k-\beta+1)^2}{8}.
\]
Indeed its finite-sum expression is
\(\ell(k-\beta-\ell+1)/2\), a concave quadratic in \(\ell\)
whose maximum over the real line is the right-hand side.

Put \(\phi(z)=6a_g(z)-a_f(z)\) and
\(\tau=k-B+1\). For fixed \(f\), the values of \(a_f\) on
\((0,1/2)\) differ by at most one. Thus the values of \(\phi\)
lie in an interval of length seven. Also
\[
 \int_0^{1/2}a_f(z)\,dz=f,\qquad
 T:=\int_0^{1/2}\phi(z)\,dz=6g-f.
\]
For a function taking values in an interval of length seven, its
variance is at most \(49/4\). To see this directly, if the interval
is \([a,a+7]\), average the nonnegative product
\((\phi-a)(a+7-\phi)\); rearrangement bounds the variance by
\((\bar\phi-a)(a+7-\bar\phi)\le49/4\).
Since our interval of integration has length \(1/2\), this gives
\[
 \int_0^{1/2}\phi(z)^2\,dz-2T^2\le\frac{49}{8}.
\]
Now
\begin{align*}
 -u\tau+\frac14\int_0^{1/2}(\tau-\phi(z))^2\,dz
 ={}&\frac18(\tau-2T-4u)^2-2uT-2u^2\\
 &+\frac14\left(\int_0^{1/2}\phi(z)^2\,dz-2T^2\right).
\end{align*}
Choose the integer \(\tau\) nearest \(2T+4u\). The square contributes
at most \(1/32\), and the variance contributes at most \(49/32\).
This choice makes \(k\) an integer nearest
\[
 2T+4u+B-1=\frac{13}{5}x-5.
\]
It lies in the admissible range \(-2q-6\le k\le8q+20\)
for \(x\ge3\).

Writing \(\eta=\{2u\}\), direct substitution in \eqref{growth:J} gives
\[
 2J(u)=2u^2-u+\frac{\eta(1-\eta)}2
 \le2u^2-u+\frac18.
\]
Insert these bounds into \eqref{growth:Ein}. The quadratic terms
cancel, and the three constant costs are
\(1/32+49/32+1/8=27/16\). Therefore
\begin{equation}\label{growth:tail-profile}
 -L_p\le p\left(x\mathcal F(x)+\frac{27}{16}\right)+O_M(1),
 \qquad
 \mathcal F(x)=4\lambda+2\lambda\{x\}
                       -12\lambda\{\alpha x\}.
\end{equation}

The remaining integral has a useful elementary bound. The following
calculation is from
\cite[\texttt{Growth/TailIntegral.lean}]{firsching}; its periodic
primitives are the ones used in Fauzan's tail analysis
\cite[Section 5.3]{fauzan}.
They are successive periodic Bernoulli primitives of the counting
error, as in Euler--Maclaurin summation. Put
\(\widetilde B_j(x)=B_j(\{x\})\), where
\[
 B_1(v)=v-\frac12,\qquad
 B_2(v)=v^2-v+\frac16,\qquad
 B_3(v)=v^3-\frac32v^2+\frac12v.
\]
Since \(6\lambda/\alpha=74\),
\[
 \mathcal F(x)+\lambda
 =2\lambda\widetilde B_1(x)
       -12\lambda\widetilde B_1(\alpha x).
\]
Define \(\overline P=(74-\lambda)/6=2923/240\) and
\begin{align*}
 P(x)-\overline P
   &=\lambda\widetilde B_2(x)-74\widetilde B_2(\alpha x),\\
 C(x)&=\frac{\lambda}{3}\widetilde B_3(x)
           -\frac{74}{3\alpha}\widetilde B_3(\alpha x).
\end{align*}
The identities \(B_2'=2B_1\), \(B_3'=3B_2\)
\cite[24.4]{dlmf} give, between breakpoints,
\[
 P'(x)=\mathcal F(x)+\lambda,\qquad
 C'(x)=P(x)-\overline P.
\]
All terms have common period \(40\). The Bernoulli polynomials have
zero mean, explaining the constants removed at each integration.
Moreover \(B_2\) and \(B_3\) have matching endpoint values, so \(P,C\)
are continuous and piecewise integration by parts introduces no jump
terms. For the bounds below it is useful to write
\[
 G_0(v)=-\frac{B_3(v)}3=\frac{v(1-v)(2v-1)}6,\qquad
 C(x)=\frac{74}{\alpha}G_0(\{\alpha x\})-\lambda G_0(\{x\}).
\]

For explicit bookkeeping put
\[
 \Psi(x)=\frac{P(x)}{x^2}+\frac{2C(x)}{x^3}
       -\frac{\overline P}{x^2}+\frac{\lambda}{x}
       -\frac{27}{32x^2}.
\]
Its derivative between breakpoints is
\[
 \Psi'(x)=\frac{x\mathcal F(x)+27/16}{x^3}
                -\frac{6C(x)}{x^4}.
\]
On \(0\le v\le1\), \(|G_0(v)|\le1/60\). For example, setting
\(w=v(1-v)\), the identity
\[
 w^2(1-4w)-\frac1{108}
 =-\left(w-\frac16\right)^2\left(4w+\frac13\right)
\]
bounds \(36G_0(v)^2\) by \(1/108<1/100\).
Consequently \(|C(x)|\le(74/\alpha+\lambda)/60<17\).
It follows that
\begin{align}
 \int_{20}^{400}\frac{x\mathcal F(x)+27/16}{x^3}\,dx
 &\le\Psi(400)-\Psi(20)
             +34\left(\frac1{20^3}-\frac1{400^3}\right)\notag\\
 &=-\frac{6843153}{128000000}.
 \label{growth:tail-integral}
\end{align}
For the last equality one only needs
\[
 P(20)=\frac{37}{2},\quad C(20)=0,\quad
 P(400)=C(400)=0.
\]
These endpoint substitutions are also checked exactly by the
accompanying script.

\subsection{The outer profile and its rank loss}

For \(K/3<p\le2h\), we have \(20/37\le x<3\), \(q=\lfloor x\rfloor\),
and \(\lfloor\alpha x\rfloor=0\). For a nonzero residue class with
\(a=a_f(c/p)\), \(d=a_g(c/p)\), the sum of its outer weights is
\begin{equation}\label{growth:GO}
 G_q(a,d)=
 \sum_{i=0}^{(2q+a-2)_+-1}
 \min\left(0,i+3d-\frac{2q+a+4}{2}\right).
\end{equation}
An empty sum is zero. This is the local weight formula after counting
the indices in the class: \(2q+a\) indices up to \(K\), \(d\) up to
\(N\), and \((2q+a-2)_+\) above \(p\). The last count determines
which rows can have a nonzero negative weight. The zero class
contributes a bounded amount, since it contains at most two pole
indices in this range.
\cite[\texttt{Growth/OuterPrime.lean}, \texttt{sum\_wO\_class}]{firsching}

The full rank loss is
\[
 r_p=\max(0,52n+2-2p)
    =p\max\left(0,\frac{13}{10}x-2\right)+O(1).
\]
Using the same factorial and class-count estimates as before gives
\begin{align}
 E_{\rm out}(x)
 ={}&-2\lambda x\lfloor x\rfloor+2J(\lambda x)
 -2\int_0^{1/2}G_{\lfloor x\rfloor}(a_f(z),a_g(z))\,dz\notag\\
 &+\max\left(0,\frac{13}{10}x-2\right),\label{growth:Eout}\\
 -L_p\le{}&pE_{\rm out}(K/p)+O(1).\label{growth:outer-profile-bound}
\end{align}
The last term is retained in full. This is the revised profile of
\cite[\texttt{Growth/OuterPrime.lean}, \texttt{outer\_Lp\_le}]{firsching}.

The profile is affine on finitely many intervals. The following table
gives \(E_{\rm out}(x)=Ax+B\); neighboring intervals with the same
formula have been combined. Endpoint values are immaterial for the
prime-sum limit.
\begin{center}
\begin{tabular}{c c r r}
\hline
Left endpoint&Right endpoint&\(A\)&\(B\)\\ \hline
\(20/37\)&\(1\)&\(37/20\)&\(-1\)\\
\(1\)&\(40/37\)&\(1\)&\(-2\)\\
\(40/37\)&\(20/13\)&\(42/5\)&\(-10\)\\
\(20/13\)&\(60/37\)&\(97/10\)&\(-12\)\\
\(60/37\)&\(80/43\)&\(231/20\)&\(-15\)\\
\(80/43\)&\(2\)&\(51/10\)&\(-3\)\\
\(2\)&\(80/37\)&\(17/4\)&\(-5\)\\
\(80/37\)&\(100/37\)&\(429/40\)&\(-19\)\\
\(100/37\)&\(120/43\)&\(503/40\)&\(-24\)\\
\(120/43\)&\(3\)&\(36/5\)&\(-9\)\\ \hline
\end{tabular}
\end{center}
To reproduce the table, use the partition \eqref{growth:grid}, adding
the rank-loss breakpoint \(20/13\), and substitute the fixed class
counts into \eqref{growth:GO}--\eqref{growth:Eout}.
Integrating the ten rows with \eqref{growth:affine-integral} gives
\begin{equation}\label{growth:outer-integral}
 \int_{20/37}^{3}\frac{E_{\rm out}(x)}{x^3}\,dx
 =\frac{129101}{96000}
 =\frac{127751}{96000}+\frac9{640}.
\end{equation}
The extra \(9/640\) is the cost of the revised full rank loss relative
to the original outer integral. The table and this equality are checked
in \cite[\texttt{Growth/OuterTable.lean},
\texttt{TableIntegrals.lean}]{firsching} and independently reproduced
by the accompanying rational-arithmetic script.

\subsection{The arithmetic budget}

Apply Lemma~\ref{growth:prime-sum} to the finite inner range, the inner
tail, and the outer range. Combine
\eqref{growth:small}, \eqref{growth:inner-integral},
\eqref{growth:tail-integral}, and \eqref{growth:outer-integral}.
Every bounded per-prime error contributes \(o(K^2)\), as shown above.
We obtain
\begin{equation}\label{growth:limsup}
 \limsup_{\substack{K\to\infty\\40\mid K}}
 \frac{\log m_K}{K^2}\le B,
\end{equation}
where the four contributions are
\[
\begin{array}{c|c}
\text{Prime range}&\text{Upper bound for the coefficient of }K^2\\ \hline
p\le K/400&111/8000\\
K/400<p\le K/20&-6843153/128000000\\
K/20<p\le K/3&322437603634266857629/7535670527041937280000\\
K/3<p\le2h&129101/96000
\end{array}
\]
and their exact sum is
\begin{equation}\label{growth:B}
 B=\frac{2031621487545722214389669}{1507134105408387456000000}
  =1.348003127429204\ldots<1.35.
\end{equation}
In particular, \eqref{growth:integrality} and the strict inequality in
\eqref{growth:B} imply, for all sufficiently large admissible \(K\),
\begin{equation}\label{growth:final}
 \boxed{\quad m_KF_K\in\mathbb Z[X],\qquad
               m_K>0,\qquad \log m_K\le1.35K^2.\quad}
\end{equation}
Indeed, choose any positive \(\varepsilon<1.35-B\) in the definition
of the upper limit.

The source formalization rounds its arithmetic bound more generously
to \(1.36+\varepsilon\). We have retained its four explicit
contributions and checked their exact sum; this justifies the sharper
eventual rounding used in \eqref{growth:final}. The difference is
only a choice of numerical margin, not a new local divisibility result.
